% part: intuitionistic-logic
% chapter: propositions-as-types

\documentclass[../../../include/open-logic-chapter]{subfiles}

\begin{document}

\olchapter{int}{pty}{రకాలుగా ప్రతిపాదనలు}

\begin{editorial}
  కర్రీ--హోవర్డ్ అనురూపతపై ఈ అధ్యాయం \emph{అత్యంత ప్రయోగాత్మకమైన} ముసాయిదా. మరింత వివరణ, ప్రేరణ అవసరం; పొరపాట్లు, లోపాలు ఉండే అవకాశం ఉంది. సాధారణీకరణ నిరూపణను సమీక్షించి విస్తరించాలి. లబ్ధ రకానికి ఉదాహరణలు లేవు. క్రమమార్పు, సరళీకరణ పరివర్తనలను చర్చించలేదు. ఇక్కడ ప్రాథమికంగా ముందే తెలిసినదిగా తీసుకున్న (రకాలతో కూడిన) లాంబ్డా కలనశాస్త్రంపై పాఠ్యం కూడా చేరితే ఇది మరింత అర్థవంతమవుతుంది. అత్యంత జాగ్రత్తగా వాడండి.
\end{editorial}

\olimport{introduction}
%\olimport{sequent-natural-deduction}
\olimport{proof-terms}
\olimport{proofs-to-terms}
\olimport{terms-to-proofs}
\olimport{types}
\olimport{reduction}
\olimport{type-preservation}
\olimport{normalization}

\OLEndChapterHook

\end{document}
